The packaging of high energy physics software with robust, yet flexible, distribution methods is a complicated problem that has been met with multiple approaches by the community.
The HEP Packaging Coordination community project expands packaging of the HEP software ecosystem through building and distributing language-agnostic conda packages on the conda-forge package index.
Through use of the conda-forge community build cyberinfrastructure, computing platform specific optimized builds of packages can be created for selections of Linux, macOS, and Windows across x86-64, AArch64/ARM64, and ppc64le architectures.
In addition to supporting builds of ROOT, this work provides multi-platform packaging of a wide array of low-level-language phenomenology tools, the broader simulation stack, end-user-analysis tools and statistical frameworks, and the reinterpretation ecosystem.
Ongoing work is also supporting builds of LHCb experiment software and distributions of community software with experiment-specific patches applied for use in LHC physics analyses.

This process significantly lowers technical barriers across tool development by providing automatic packaging systems with source code, distribution through secure and transparent build cyberinfrastructure, and enables use through multi-platform optimized binary builds.
When combined with next generation scientific package management and manifest tools, the creation of fully specified, portable, and trivially reproducible multi-language software environments becomes easy and fast, even with the use of development platforms for hardware accelerators (e.g. CUDA on NVIDIA GPUs).
This talk provides an overview of the work, gives practical recommendations for adoption and best practices for both software maintainers and end-user analysts, and demonstrates examples of new distribution methods that are complementary to existing community technologies, such as CernVM-FS.
